% OWNER: methods
% STATUS: draft
% SOURCES: GLOSSARY.md §5; HIPOTEZY_REWIZJA.md; CoreProfileWizard / survey flows (opis parafrazowany)
\chapter{Treść instrumentów i ekranów}
\label{app:b:top}
\ChapterStatus{draft}{methods}

Załącznik uzupełnia tabelę master z~\cref{tab:ch06:instruments}. Teksty poniżej są \textbf{parafrazami / skrótami} na potrzeby dokumentacji metody --- nie stanowią pełnej reprodukcji pozycji chronionych licencją instrumentów (w szczególności IPIP). Ekrany UI opisano funkcjonalnie, bez zrzutów zawierających dane uczestników.

\section{PRS --- mood grid}
\label{app:b:sec:prs}

\begin{itemize}
  \item \textbf{Cel:} uchwycenie postrzeganej regeneracyjności / nastroju przestrzeni na uproszczonej płaszczyźnie (adaptacja kategorii znanych ze skali PRS; \cite{hartig1997prs}).
  \item \textbf{Forma:} interaktywna siatka (mood grid) zamiast pełnej baterii itemów papierowych.
  \item \textbf{Pomiary w IDA:} \texttt{prs\_ideal\_*} (profil), \texttt{prs\_current\_*} i \texttt{prs\_target\_*} (Room Setup).
  \item \textbf{Czego nie ma:} pomiaru \emph{post} po wyborze wizualizacji AI.
\end{itemize}

\section{Tinder / IAT-like}
\label{app:b:sec:tinder}

\begin{itemize}
  \item \textbf{Cel:} preferencje ukryte wobec stylów / materiałów / nastrojów wizualnych poprzez szybkie decyzje binarne (inspiracja logiką IAT; \cite{greenwald1998iat}).
  \item \textbf{Forma:} gest swipe (akceptacja / odrzucenie) na zestawie bodźców obrazowych.
  \item \textbf{Zapis:} \texttt{participant\_swipes} (m.in.\ kierunek, \texttt{reaction\_time\_ms}); agregaty \texttt{implicit\_*} w \texttt{participants}.
  \item \textbf{Limit:} brak osobnego pomiaru dwell/hesitation w bazie.
\end{itemize}

\section{Preferencje jawne (semantic differential i pokrewne)}
\label{app:b:sec:explicit}

\begin{itemize}
  \item Slidery / skale (ciepło, jasność, złożoność, faktura), wybór palety i materiałów.
  \item Zmienne: \texttt{explicit\_warmth}, \texttt{explicit\_brightness}, \texttt{explicit\_complexity}, \texttt{explicit\_palette}, \texttt{explicit\_material\_*} itd.
  \item Po profilu system może zapisać porównanie z implicit (\texttt{preference\_comparison\_json}).
\end{itemize}

\section{Biophilia i inspiracje}
\label{app:b:sec:biophilia}

\begin{itemize}
  \item Wybór dawki biophilii (uproszczone poziomy) $\rightarrow$ \texttt{biophilia\_score}.
  \item Upload / wybór inspiracji (1--$n$) $\rightarrow$ tagi i agregaty \texttt{inspiration\_*}; wspierają źródło \texttt{inspiration\_reference}.
\end{itemize}

\section{IPIP-NEO-120 (Big Five)}
\label{app:b:sec:ipip}

\begin{itemize}
  \item \textbf{Instrument:} IPIP-NEO-120 (120 pozycji; 5 domen + 30 facetów). Rama taksonomiczna Big Five: \cite{john2008bigfive}.
  \item \textbf{Forma w IDA:} dedykowany krok \texttt{/flow/big-five} (ścieżka Full).
  \item \textbf{Zapis:} domeny \texttt{big5\_*}, \texttt{big5\_facets}, surowiec \texttt{big5\_responses} (JSONB).
  \item \textbf{Licencja / reprodukcja:} w rozprawie \emph{nie} publikuje się pełnej listy itemów IPIP; dokumentuje się wersję instrumentu i sposób scoringu. Przykładowa parafraza formatu odpowiedzi: skala Likerta zgody z twierdzeniem o zachowaniu / preferencji (treść itemu wg oficjalnego IPIP).
\end{itemize}

\section{Room Setup i laddering}
\label{app:b:sec:room}

\begin{itemize}
  \item Zdjęcie pomieszczenia + kontekst użycia (solo/shared), aktywności, pain points.
  \item Laddering: od atrybutów przestrzeni przez konsekwencje do potrzeby rdzeniowej (\texttt{ladder\_path}, \texttt{ladder\_core\_need}).
  \item Wspiera źródło macierzy \texttt{mixed\_functional}.
\end{itemize}

\section{Ślepa macierz sześciu źródeł}
\label{app:b:sec:macierz}

Uczestnik widzi sześć wizualizacji bez etykiet. Po wyborze system zapisuje źródło odpowiadające wybranej komórce. Źródła (kanon kodu):

\begin{enumerate}
  \item \texttt{implicit} --- wariant preferencji ukrytych,
  \item \texttt{explicit} --- wariant preferencji jawnych,
  \item \texttt{personality} --- wariant osobowościowy,
  \item \texttt{mixed} --- wariant mieszany,
  \item \texttt{mixed\_functional} --- wariant mieszany z funkcją przestrzeni,
  \item \texttt{inspiration\_reference} --- wariant referencji inspiracji.
\end{enumerate}

\section{Ankiety UX po flow}
\label{app:b:sec:ux}

\begin{itemize}
  \item \textbf{SUS:} bateria stwierdzeń o użyteczności (klucze \texttt{sus\_1}\ldots\texttt{sus\_10} w UI); agregat \texttt{sus\_score}.
  \item \textbf{Agency / clarity / satisfaction:} krótkie skale własne (poczucie sprawczości, jasność procesu, satysfakcja); agregaty \texttt{agency\_score}, \texttt{clarity\_score}, \texttt{satisfaction\_score}.
  \item Satysfakcja jest mierzona na poziomie uczestnika, nie per wybrany \texttt{selected\_source}.
\end{itemize}
